\documentclass[11pt]{article}
\usepackage[top=1.00in, bottom=1.0in, left=1in, right=1in]{geometry}
\renewcommand{\baselinestretch}{1.1}
\usepackage{graphicx}
\usepackage{natbib}
\usepackage{amsmath}
\usepackage{parskip}
\usepackage{hyperref}
\usepackage[symbol]{footmisc} % https://tex.stackexchange.com/questions/826/symbols-instead-of-numbers-as-footnote-markers

\def\labelitemi{--}
\parindent=0pt

\begin{document}
\bibliographystyle{/Users/Lizzie/Documents/EndnoteRelated/Bibtex/styles/besjournals}
\renewcommand{\refname}{\CHead{}}

% \setlength{\parindent}{0cm}
% \setlength{\parskip}{5pt}

{\Large My guidelines for use of generative AI when working with me}\footnote[1]{I (Lizzie) try to share this policy with everyone, regardless of whether I think you use genAI or not.}


\emph{Last updated:} \today

Generative AI is a tool that uses large language models and their kin (e.g., `reasoning models') to perform a number of tasks related to generating text. {\bf If you are collaborating with me in any way (this includes if you are a trainee\footnote[2]{Trainee means you are an undergraduate student, graduate student or postdoc that I work with. You could be in my lab, or visiting my lab, or I could be on your committee.}) please disclose all use of generativeAI to me.} How you do this is up to you, but it's important to me that you do it in such a way that I don't end up thinking I am communicating with you directly when really there is generativeAI in between us somehow (if this ends up happening, I may ask you to follow more specific guidelines of disclosure). I also follow this guideline and will let you know when I have or am using genAI. 

My goal in asking for this is open communication, which for me happens person to person. My hope is that following this request is easy to do and generally leads to conversations about how to best develop skills, techniques, and do science (which may include how to best use LLMs). I definitely do not know the best way to do everything (though I do have lots of opinions, and I have thought hard about how to do certain things), so this should be a conversation. 


\subsection*{Writing with genAI}

While I appreciate the value of genAI in some regards, I consider writing as a scientist a very important skill. Learning to write well helps you learn to think logically, communicate well and connect with different audiences. I figure a lot of things out while being forced to write a paper, grant or thoughtful note on a problem/question (and this is often when I have to do more brainstorming and logical thinking). I don't believe there is good evidence that using generative AI (e.g., chatGPT, current versions of Grammerly etc.) to write helps you learn to write. 

If you are a trainee, I think it makes it harder for me to train you to the best of my ability if you use these products versus share your writing with me. I think it is also important to recognize and stay informed of current \href{http://grad.ubc.ca/current-students/student-responsibilities/use-generative-ai}{UBC guidelines} (and the guidelines wherever you are)\footnote[3]{At many universities failure to disclose the use of genAI is considered academic misconduct.}
 on the use of generative AI for producing and editing text. In particular, you should understand when, why and how using generative AI to generate text is plagiarism.
 
{\bf My current policy if you are in the lab and/or collaborating with me} is that I ask you to not use generative AI for any writing you will submit to me.\footnote[4]{The guidelines in this section of this file apply to your writing of text, not your writing of code, but the next section applies to your code.} I am happy to discuss this policy and my reasons for it. It's not a policy that I consider set in stone forever, but I do feel it is the best policy for now for me to be most helpful as someone training others in writing and to meet my own personal standards of writing. 

If you are collaborating with others outside the lab, I encourage you to try to discuss what use of genAI you're all okay with early on (you're welcome to share this document, and I am here to try to help you with these conversations as needed). As a first author on a manuscript, you may be responsible for a genAI statement when finalizing your manuscript, so you will need to know how your co-authors used genAI. 

This policy applies to writing of text for proposals, manuscripts, grants, emails, coding issues (on GitHub and other forums) etc..

{\bf If you have asked me to be on your supervisory committee} my rules are:
\begin{enumerate}
\item If you want to use generative AI to help you generate ideas, text or structure/refine your writing. I am not currently a good choice for your supervisory committee. 
\item I’d prefer you not use generative AI to edit your writing. I realize that means I will read text with some grammatical errors and that is okay. 
\item I understand that you may want to use it to edit your grammar.\footnote[5]{I realize genAI is increasingly embedded in many programs and I am generally okay with use for fixing typos, verb-noun agreement and extremely minor fixes of the sort that were once done without genAI.}  If you want to use generative AI as a tool to help you edit your writing, you must agree that you will not use generative AI beyond this without telling me, and show me that you know the different levels of using generative AI and how to properly disclose them. Thus, you need to:
\begin{enumerate}
\item Make up a short (say, 1-3 pages) document showing examples of both how you use AI to edit your writing \emph{and} examples where using it without quotation marks would be considered plagiarism. Please also include examples where it is editing your language and not your grammar, and confirm you will not do this in writing you submit to me.
\item You also need to find a way to share or record all of the changes made so you can document them.
\item Acknowledge that you know that using generative AI for writing that I am reviewing without disclosing it to me is \emph{academic misconduct} to me. (So please do not do this.)
\end{enumerate}
\end{enumerate}

\subsection*{About coding and other uses of generative AI...} 

{\bf If you use genAI to help you write code, you must add a comment to your code file whenever you copy and paste in code from genAI and briefly explain how you used it.} Examples include `f(x) copied from Claude [model name and version]' or `I wrote original code, but used PurpleUnicorn to update the parameters.' I recommend you make sure it's clear where the pasted code starts and stops, and include the model and version name somewhere as you may be asked for it in journal submission. % Added per 5 August 2026 conversation with Mao, Victor and Christophe.  

I consider writing code different from writing other text for several reasons, including that I am not asked to read your code line by line (instead I evaluate the output) and that you can (and should) test your code in ways you cannot test your writing. That doesn't mean I endorse using generative AI for coding. In my experience, it is not a great way to learn to code unless you have other structured ways to learn supplement coding help from generative AI. And you should invest more in learning unit testing and other ways to test code, especially if you are using generative AI heavily for coding. 

Relatedly, I don't think generative AI is necessarily a great tool to develop your ideas, hypotheses etc. and you should think hard about how you use genAI for this. 

As in all of graduate school, I encourage you to think about the best ways to learn and structure your learning so that you build skills (not just get the product and/or degree). Uncritical adoption of anything is not a very good approach, especially in science; find ways to make yourself aware of the pros and cons, and stay informed. 


\subsection*{Failure to follow these guidelines}
If you have (or think maybe you have) failed to follow these guidelines, please tell me and we can talk about it. This is a much better approach than me finding out some other way. % \footnote{If you don't tell me, there may be other consequences.}
% If you fail to follow these guidelines and we continue to work together, I may ask you to include a statement of how you used generative AI in every communication with me. That includes every email you send, all documents and all work (GitHub issues etc.). {\bf Include this statement on ALL communication, even if you did not use genAI.} If you did use it, please state exactly how, which program and provide a shareable link that shows how you used it in this instance. 


\subsection*{Examples of AI use statements} 
More and more we'll all likely need to write genAI use statements so here I am keeping some I think are especially good. 

From \href{https://arxiv.org/pdf/2607.17397}{`The unintended consequences of large language models as a labor-augmenting technology in science'} :
\begin{quote}
The text of this paper was written without AI assistance beyond the spell-checking and grammar-checking features built into Overleaf. The figures were drafted with pen and paper and hand-coded in Mathematica, but CTB used ChatGPT 5.5 assistance to refine fonts, label positions, etc., as well as for the help with the intricacies of \LaTeX formatting. Literature search was predominantly conducted using traditional tools such as Google Scholar, Google search, Semantic Scholar, and Web of Science but such tools now integrate AI. We wrote all proofs, but CTB and KG worked with ChatGPT 5.4, ChatGPT 5.5, and Gemini 3.1 to suggest proof strategies, and ED worked with Claude Fable 5 to check the proofs for mathematical and notational consistency and mistakes. MJC did not use LLMs in their contributions to the paper.
\end{quote}



\end{document}

